\paragraph{AQ-136-COTENSOR: a simplicial path theory over any coefficient ring.}
For any $C\in\mathcal A_k$ over any commutative ring $k$, put
\[
 P_s(C)_n^{w,r}=C^{w,r}(\Delta[n]\times\Delta[1]).             \tag{136.25}
\]
Here $C(X)$ is the compatible-section extension from simplices
used in AQ-136-COMPARISON.  The operations are taken pointwise.
Then $P_s(C)$ belongs to $\mathcal A_k$.  Its constant-path
inclusion $j:C\to P_s(C)$ and two evaluations
$e_0,e_1:P_s(C)\to C$ are strict PD morphisms over $S$,
and $e_0j=e_1j=1$.

\begin{proof}
The construction is a limit of the given bigraded differential
PD algebras, so inherits every operation and the coefficient
map.  Functoriality in $\Delta[n]$ gives its simplicial
structure; projection and the two vertex inclusions give
the displayed maps and their identities.

Fix a cochain bound $d$.  AQ-136-COMPARISON gives a weight
bound $W_d$ independent of the simplicial set.  Each
$\Delta[n]\times\Delta[1]$ is contractible: it is the nerve
of the product ordered set $[n]\times[1]$, which has a
terminal element, or use the usual monotone contraction of
each simplex.  Therefore, for $w\geq W_d$,
\[
 H^0(P_s(C)_n^{w,*})=S^w,\qquad
 H^r(P_s(C)_n^{w,*})=0\quad(1\leq r\leq d),                \tag{136.26}
\]
with the first isomorphism induced by the coefficient map.
The bound is uniform in $n$, as required by (136.3).
There is no interchange of infinitely many separate
simplex-dependent bounds.

For fixed $(w,r)$, write $F=C_*^{w,r}$.  The corresponding
simplicial module in (136.25) is its cotensor
$F^{\Delta[1]}$.  Evaluation at the zero vertex and the
constant map exhibit a simplicial linear homotopy
equivalence
\[
                       F^{\Delta[1]}\simeq F.              \tag{136.27}
\]
Explicitly the order-preserving minimum map
$[1]\times[1]\to[1]$ induces a simplicial map
$\mu:\Delta[1]\times\Delta[1]\to\Delta[1]$
with $\mu(s,0)=0$ and $\mu(s,1)=s$.
Precomposition by this map gives a homotopy between the
constant-path evaluation and the identity on the cotensor.
All its maps are linear, since operations on functions to
$F$ are pointwise.  Thus it induces the same normalized
simplicial homology as $F$, without a flatness hypothesis.
Condition (136.4) for $C$ consequently gives that condition
for $P_s(C)$, with the same cutoff function.  This proves
admissibility and all assertions.
\end{proof}

\paragraph{AQ-136-PLOCAL-LOCALIZATION: the $p_0$-local groupoid is terminal.}
For every commutative $\mathbb Z_{(p_0)}$-algebra $k$, the
ordinary localization of $\mathcal A_k$ at all its morphisms
is equivalent to the terminal groupoid.  The same holds
for a finite product of such coefficient rings and hence
for every ring of positive characteristic.  Together with
AQ-136-PLOCAL this proves both non-strict original
uniqueness assertions in these coefficient ranges.

\begin{proof}
We first show that any two parallel morphisms
$f,g:C\rightrightarrows D$ have the same image in the
all-arrow localization.  Choose $N=p_0^e$ so large that
both $C_*^{N,0}$ and $D_*^{N,0}$ are acyclic.  The proof
of AQ-136-PLOCAL provides an admissible source and map
\[
                         a:B_N^{\rm PD}\longrightarrow C.
\]
Set $F=D_*^{N,0}$.  The coordinate images under $fa$ and
$ga$ define simplicial linear maps
$\phi_0,\phi_1:L\to F$.  On the constant submodule
$kt\subset L$ both send $t=\sum_jx_j$ to the coefficient
$T^{[N]}$.  Their difference thus factors through
\[
             Q=L/kt,\qquad \psi:Q\longrightarrow F.        \tag{136.28}
\]
Normalization is exact, and the disk calculation for $L$
shows that $N(Q)=k[1]$.
The normalized map $N(\psi)$ is consequently a cycle in
$N(F)_1$.  Since $H_1(N(F))=0$, that cycle is the boundary
of an element of $N(F)_2$.  Mapping the generator of
$k[1]$ to this element gives a chain homotopy between
$N(\psi)$ and zero.  The chain-to-simplicial homotopy
statement, Stacks Project Lemma14.29.3/Tag\texttt{01A4},
gives a simplicial linear homotopy from zero to $\psi$.
Equivalently it gives a linear map
$Q\to F^{\Delta[1]}$ with those endpoints.

Compose this map with $L\to Q$ and add the constant-path
map associated with $\phi_0$.  The result is a simplicial
linear map
\[
 \Phi:L\longrightarrow F^{\Delta[1]},\qquad
 e_0\Phi=\phi_0,\quad e_1\Phi=\phi_1,\quad
 \Phi(t)=T^{[N]}\text{ as a constant path}.                 \tag{136.29}
\]
The last equality is exact: the homotopy term factors
through $Q$ and is zero on $kt$, throughout the homotopy.
This relative condition is what preserves the original
coefficient relation.

Apply the free PD polynomial universal property in each
simplex, sending $x_j$ to $\Phi(x_j)$, $e_j$ to
$d\Phi(x_j)$, and every coefficient of $S$ to itself.
Equation (136.29) gives the two relations
$\sum x_j=T^{[N]}$ and $\sum e_j=0$; all operators are
natural in the simplex.  We obtain a strict differential
PD map over $S$
\[
 H:B_N^{\rm PD}\longrightarrow P_s(D),\qquad
                         e_0H=fa,\quad e_1H=ga.            \tag{136.30}
\]
It preserves both gradings, because the coordinate
homotopy was carried out within weight $N$, cochain
degree zero.  AQ-136-COTENSOR proves that its target is
admissible over the same coefficient ring.

Since $e_0j=e_1j=1$, the two evaluations both equal
$j^{-1}$ after inversion of all morphisms.  Thus
$fa=ga$ in the localization, and inversion of $a$
gives $f=g$ there.
Finite common sources exist by AQ-136-PLOCAL.
For any two finite zigzags with the same endpoints,
choose one common source for all their objects.  For
an original arrow $u:C_i\to C_j$, its composite with
the chosen source arrow to $C_i$ and the chosen source
arrow to $C_j$ are parallel; they have the same localized
image by the argument just given.  Every edge is
therefore the corresponding ratio of these source
arrows, and both zigzags telescope to the same ratio.
The localization is nonempty and connected, and has
exactly one arrow between any two objects.

For a finite product of coefficient rings, choose the
sources $B_{N_i}^{\rm PD}$ and the relative coordinate
homotopies separately in each factor.  Their finite
products are admissible and form the same diagrams over
the product ring: finite idempotent decomposition
commutes with compatible sections and hence with
$P_s$.  The identical parallel-arrow and telescope
argument applies.  The Chinese-remainder decomposition
proved in AQ-136-PLOCAL covers every positive-characteristic
ring.  The zero ring again causes no exception.
\end{proof}

The simplicial cotensor varies the argument $X$ of a
theory.  It does not adjoin a weight-zero element with
endpoints zero and one to one cochain algebra; the
hypotheses of AQ-136-PATH-OBSTRUCTION are therefore
not being asserted.  All coefficient divided powers
remain fixed in (136.29)--(136.30).
The two non-strict questions over $\mathbb Z$ remain
unresolved by this argument, since its admissible
source $B_N^{\rm PD}$ uses the one-prime free coefficient
decomposition (136.23).

\paragraph{AQ-136-WEIGHT-ONE: admissibility does not force the lowest row to be acyclic.}
For $N\geq2$ and $k\ne0$, there is an admissible theory
$C$ to which the standard source $B=B_1$ has no
$S$-morphism.  Define, inside (136.5),
\[
 C^{w,r}_n=
 \begin{cases}
 S^w,&w<N,\ r=0,\\
 0,&w<N,\ r>0,\\
 B^{w,r}_n,&w\geq N.
 \end{cases}                                               \tag{136.31}
\]

\begin{proof}
These submodules are preserved by $d$ and simplicial
operators.  Products involving a term of weight at
least $N$ remain in that range; products of lower-weight
coefficients are coefficients.  Divided powers likewise
preserve the subalgebra: they preserve the coefficient
algebra, and multiply a positive weight at least $N$
by a positive integer.  Thus $C$ is a strict differential
PD $S$-subalgebra of $B$.
Its vertical cohomology is $S$ in degree zero and zero
otherwise in every weight.  All sufficiently large
weight rows are the rows of $B$, so (136.4) holds.
Consequently $C$ is admissible.

Its weight-one zero-form row is the nonzero constant
simplicial module $kT$.  A hypothetical $S$-map
$B\to C$ would restrict on the coordinate module to
a simplicial linear map $L\to\underline{kT}$ taking
$t$ to $T$.  On normalization, a chain map from the
disk $k\xrightarrow{\pm1}k$ to $k[0]$ must have zero
degree-zero component, because its composite with
the disk differential is zero.  It cannot take $t$
to nonzero $T$.  This is the required contradiction.
\end{proof}
